%═══════════════════════════════════════════════════════════════════════════════
\chapter{Algoritmos de referencia}
\label{ap:algoritmos}
%═══════════════════════════════════════════════════════════════════════════════

Este apéndice consolida los algoritmos que los ejercicios computacionales del
libro requieren implementar. No repite los algoritmos que ya aparecen completos
en el cuerpo del texto (Algoritmos~\ref{alg:ciclo}, \ref{alg:adj_ode},
\ref{alg:sinkhorn}, \ref{alg:ddpm}, \ref{alg:rectified_flow}), sino los que se
mencionan en los ejercicios sin estar escritos explícitamente.

Para cada algoritmo se proporciona: pseudocódigo autocontenido, los parámetros
con sus valores por defecto recomendados, y una nota sobre las trampas numéricas
más frecuentes en la implementación.

%───────────────────────────────────────────────────────────────────────────────
\section{Optimizadores}
\label{sec:ap_optimizadores}
%───────────────────────────────────────────────────────────────────────────────

\subsection{Adam}
\label{subsec:ap_adam}

\textit{Definido en: Capítulo~\ref{ch:cap5} (Definición~\ref{def:adam}).
Usado en: todos los ejercicios computacionales con entrenamiento de redes.}

\medskip
Adam combina momento de primer orden (media exponencial del gradiente) y
segundo orden (media exponencial del cuadrado del gradiente) con corrección
de sesgo de arranque.

\begin{algorithm}[H]
\caption{Adam -- paso de actualización}
\label{alg:ap_adam}
\KwIn{Gradiente $\bm{g}^{(k)}$; estado $({\bm{m}}^{(k-1)}, {\bm{v}}^{(k-1)}, k)$;
      parámetros $(\eta, \beta_1, \beta_2, \varepsilon)$}
\KwOut{Actualización $\Delta\bm{\theta}^{(k)}$; nuevo estado}

$\bm{m}^{(k)} \leftarrow \beta_1\bm{m}^{(k-1)} + (1-\beta_1)\bm{g}^{(k)}$\;
$\bm{v}^{(k)} \leftarrow \beta_2\bm{v}^{(k-1)} + (1-\beta_2)(\bm{g}^{(k)})^2$\;
$\hat{\bm{m}} \leftarrow \bm{m}^{(k)}/(1-\beta_1^k)$\tcp*{corrección de sesgo}
$\hat{\bm{v}} \leftarrow \bm{v}^{(k)}/(1-\beta_2^k)$\tcp*{corrección de sesgo}
$\Delta\bm{\theta}^{(k)} \leftarrow -\eta\,\hat{\bm{m}}/(\sqrt{\hat{\bm{v}}}+\varepsilon)$\;
\Return $\Delta\bm{\theta}^{(k)}$, $(\bm{m}^{(k)}, \bm{v}^{(k)}, k+1)$\;
\end{algorithm}

\medskip
\noindent\textbf{Parámetros por defecto:}
$\beta_1 = 0.9$, $\beta_2 = 0.999$, $\varepsilon = 10^{-8}$.
La tasa $\eta$ es el único hiperparámetro que requiere búsqueda; el rango
típico es $[10^{-4}, 10^{-2}]$.

\medskip
\noindent\textbf{Trampas numéricas:}
\begin{itemize}
  \item Inicializar el estado con $\bm{m}^{(0)} = \bm{v}^{(0)} = \bm{0}$ y
        $k = 1$ antes del primer paso. Olvidar la inicialización produce un
        estado indefinido.
  \item El denominador $\sqrt{\hat{\bm{v}}} + \varepsilon$ puede volverse
        numéricamente cero si $\varepsilon$ es demasiado pequeño para la
        precisión usada. En \texttt{float32}, usar $\varepsilon \geq 10^{-7}$.
  \item El error más frecuente: llamar al paso de Adam sin incrementar $k$.
        Sin la corrección de sesgo correcta, las primeras iteraciones producen
        pasos exageradamente grandes.
\end{itemize}

\subsection{Scheduler coseno con warmup lineal}
\label{subsec:ap_scheduler}

\textit{Usado en: Capítulos~\ref{ch:cap8} (entrenamiento de MLPs en HEP),
\ref{ch:cap12} (entrenamiento de score networks).}

\medskip
El scheduler coseno decrece $\eta$ suavemente desde $\eta_{\max}$ hasta
$\eta_{\min}$ con una fase inicial de warmup lineal:

\begin{algorithm}[H]
\caption{Scheduler coseno con warmup lineal}
\label{alg:ap_scheduler}
\KwIn{Paso actual $k$; total de pasos $K$; fracción de warmup $w\in(0,1)$;
      $\eta_{\min}$, $\eta_{\max}$}
\KwOut{Tasa de aprendizaje $\eta^{(k)}$}

$k_w \leftarrow \lfloor wK \rfloor$\tcp*{pasos de warmup}
\eIf{$k \leq k_w$}{
  $\eta^{(k)} \leftarrow \eta_{\max}\cdot k/k_w$\tcp*{rampa lineal}
}{
  $t \leftarrow (k - k_w)/(K - k_w)$\;
  $\eta^{(k)} \leftarrow \eta_{\min} +
  \frac{1}{2}(\eta_{\max}-\eta_{\min})(1 + \cos(\pi t))$\tcp*{coseno}
}
\Return $\eta^{(k)}$\;
\end{algorithm}

\noindent\textbf{Parámetros típicos:} $w = 0.05$--$0.10$, $\eta_{\min} =
\eta_{\max}/100$.

%───────────────────────────────────────────────────────────────────────────────
\section{Integración numérica de SDEs}
\label{sec:ap_sdes}
%───────────────────────────────────────────────────────────────────────────────

\textit{Usado en: Capítulos~\ref{ch:cap10} (simulación del proceso OU,
estimación de parámetros), \ref{ch:cap11} (reverse SDE de Anderson),
\ref{ch:cap12} (muestreo de DDPM/DDIM).}

\subsection{Euler-Maruyama}
\label{subsec:ap_euler_maruyama}

El método de Euler-Maruyama es la discretización más simple de la SDE
$dX_t = \mu(X_t,t)\,dt + \sigma(X_t,t)\,dW_t$:

\begin{algorithm}[H]
\caption{Euler-Maruyama para una SDE escalar o vectorial}
\label{alg:ap_em}
\KwIn{Condición inicial $X_0$; coeficientes $\mu(\cdot,\cdot)$,
      $\sigma(\cdot,\cdot)$; intervalo $[0,T]$; paso $\Delta t$}
\KwOut{Trayectoria $\{X_{k\Delta t}\}_{k=0}^{T/\Delta t}$}

$X \leftarrow X_0$\;
\For{$k = 0, 1, \ldots, T/\Delta t - 1$}{
  $t \leftarrow k\Delta t$\;
  $\xi \leftarrow \mathcal{N}(\bm{0}, \bm{I})$\tcp*{ruido estándar}
  $X \leftarrow X + \mu(X,t)\,\Delta t + \sigma(X,t)\sqrt{\Delta t}\,\xi$\;
}
\Return trayectoria\;
\end{algorithm}

\noindent\textbf{Orden de convergencia:} fuerte $1/2$, débil $1$.
Esto significa: el error en la trayectoria individual (norma $L^2$ de la
diferencia) decrece como $O(\sqrt{\Delta t})$; el error en las estadísticas
(esperanza, varianza) decrece como $O(\Delta t)$.

\medskip
\noindent\textbf{Trampas numéricas:}
\begin{itemize}
  \item El incremento estocástico es $\sigma(X,t)\sqrt{\Delta t}\,\xi$ con
        $\xi\sim\mathcal{N}(0,1)$, \textbf{no} $\sigma(X,t)\Delta W$ con
        $\Delta W\sim\mathcal{N}(0,\Delta t)$: son equivalentes pero la primera
        forma es numericamente más estable al evitar muestrear con varianza
        pequeña.
  \item Para el proceso OU ($\mu(X) = -\alpha X$), la condición de estabilidad
        requiere $\Delta t < 2/\alpha$. Para $\alpha = 1$: $\Delta t < 2$. En
        la práctica usar $\Delta t \leq 0.1/\alpha$.
  \item En la vectorial, $\xi\in\mathbb{R}^d$ debe tener componentes
        \textbf{independientes}. Usar \texttt{np.random.randn(d)} y no
        \texttt{np.random.multivariate\_normal} (más lento sin ganancia).
\end{itemize}

\subsection{Milstein}
\label{subsec:ap_milstein}

Milstein añade una corrección de orden 2 en el incremento estocástico que
mejora el orden fuerte a 1:

\begin{algorithm}[H]
\caption{Método de Milstein (caso escalar)}
\label{alg:ap_milstein}
\KwIn{Igual que Euler-Maruyama}
\KwOut{Trayectoria con orden fuerte 1}

$X \leftarrow X_0$\;
\For{$k = 0, 1, \ldots$}{
  $t \leftarrow k\Delta t$\;
  $\xi \leftarrow \mathcal{N}(0,1)$\;
  $X \leftarrow X + \mu(X,t)\,\Delta t + \sigma(X,t)\sqrt{\Delta t}\,\xi$
  $\quad+ \tfrac{1}{2}\sigma(X,t)\partial_x\sigma(X,t)\,(\xi^2-1)\Delta t$\;
}
\end{algorithm}

\noindent\textbf{Cuándo usarlo:} cuando se necesita precisión en trayectorias
individuales (e.g., estimación de funcionales de la trayectoria) y el costo
de calcular $\partial_x\sigma$ es manejable. Para los ejercicios de los
Capítulos~\ref{ch:cap10}--\ref{ch:cap11}, Euler-Maruyama con $\Delta t$
suficientemente pequeño es adecuado.

\noindent\textbf{Nota:} para $\sigma$ constante (proceso OU, browniano puro),
$\partial_x\sigma = 0$ y Milstein = Euler-Maruyama. La corrección solo actúa
cuando $\sigma$ depende del estado.

\subsection{Reverse SDE de Anderson: discretización de Euler}
\label{subsec:ap_reverse_sde}

\textit{Usado en: Capítulo~\ref{ch:cap11} (Ejercicio computacional 1),
Capítulo~\ref{ch:cap12} (muestreo con score network).}

\medskip
El reverse SDE~\eqref{eq:score_sde_reverse} con score $\bm{s}_\theta(\bm{x},t)$
aprendido se discretiza hacia atrás en tiempo:

\begin{algorithm}[H]
\caption{Muestreo con reverse SDE (discretización de Euler)}
\label{alg:ap_reverse_sde}
\KwIn{Score network $\bm{s}_\theta$; forward drift $\bm{f}$ y difusión $g$;
      número de pasos $N$; horizonte $T$}
\KwOut{Muestra aproximada de $p_0$}

$\bm{x} \leftarrow \mathcal{N}(\bm{0}, \bm{I})$\tcp*{muestra del prior}
$\Delta t \leftarrow T/N$\;
\For{$k = N, N-1, \ldots, 1$}{
  $t \leftarrow k\Delta t$\;
  $\bm{d} \leftarrow [-\bm{f}(\bm{x},t) + g(t)^2\,\bm{s}_\theta(\bm{x},t)]
  \cdot\Delta t + g(t)\sqrt{\Delta t}\,\mathcal{N}(\bm{0},\bm{I})$\;
  $\bm{x} \leftarrow \bm{x} + \bm{d}$\;
}
\Return $\bm{x}$\;
\end{algorithm}

\noindent\textbf{Nota sobre el signo:} el tiempo corre hacia atrás ($k=N,\ldots,1$)
y el drift del reverse SDE tiene el signo opuesto al del forward más el término
de score. El error más frecuente es confundir el signo del drift o el orden de
los pasos.

\noindent\textbf{Para el proceso OU} ($\bm{f}(\bm{x}) = -\bm{x}$, $g = \sqrt{2}$):
el drift del reverse es $\bm{x} + 2\bm{s}_\theta(\bm{x},t)$.

%───────────────────────────────────────────────────────────────────────────────
\section{Noise schedules para DDPM}
\label{sec:ap_schedules}
%───────────────────────────────────────────────────────────────────────────────

\textit{Usado en: Capítulo~\ref{ch:cap12} (Ejercicios computacionales 1 y 2).}

\medskip
Los schedules determinan $\{\beta_t\}_{t=1}^T$ y por tanto
$\bar{\alpha}_t = \prod_{s=1}^t(1-\beta_s)$.

\begin{algorithm}[H]
\caption{Noise schedules lineales y coseno}
\label{alg:ap_schedule}
\KwIn{$T$ (número de pasos); tipo: \texttt{lineal} o \texttt{coseno};
      $\beta_{\min}=10^{-4}$, $\beta_{\max}=0.02$, $s=0.008$}
\KwOut{Arrays $\{\beta_t\}$, $\{\alpha_t\}$, $\{\bar\alpha_t\}$,
       $\{\sqrt{\bar\alpha_t}\}$, $\{\sqrt{1-\bar\alpha_t}\}$}

\If{\texttt{tipo == 'lineal'}}{
  $\beta_t \leftarrow \beta_{\min} + (t-1)/(T-1)\cdot(\beta_{\max}-\beta_{\min})$
  para $t=1,\ldots,T$\;
}
\If{\texttt{tipo == 'coseno'}}{
  $f(t) \leftarrow \cos^2\!\left(\frac{t/T + s}{1+s}\cdot\frac{\pi}{2}\right)$\;
  $\bar\alpha_t \leftarrow f(t)/f(0)$ para $t=0,1,\ldots,T$\;
  $\beta_t \leftarrow 1 - \bar\alpha_t/\bar\alpha_{t-1}$,
  clippeado a $[0,\, 0.999]$\;
}
$\alpha_t \leftarrow 1 - \beta_t$;
$\bar\alpha_t \leftarrow \prod_{s=1}^t\alpha_s$\tcp*{producto acumulado}
\Return arrays derivados\;
\end{algorithm}

\noindent\textbf{Nota importante:} calcular $\bar\alpha_t$ como producto
acumulado directo en \texttt{float32} produce underflow numérico para
$T\gtrsim 200$ con el schedule lineal. Usar en su lugar:
$\log\bar\alpha_t = \sum_{s=1}^t\log\alpha_s$, y exponenciar solo cuando
sea necesario.

%───────────────────────────────────────────────────────────────────────────────
\section{Métricas de evaluación}
\label{sec:ap_metricas}
%───────────────────────────────────────────────────────────────────────────────

\textit{Usado en: múltiples ejercicios computacionales para comparar
distribuciones empíricas.}

\subsection{Distancia de Wasserstein $W_1$ empírica (1D)}
\label{subsec:ap_wasserstein}

Para muestras 1D $\{x_i\}_{i=1}^N$ de $p$ y $\{y_j\}_{j=1}^M$ de $q$,
la distancia $W_1$ se calcula eficientemente por la fórmula de los cuantiles:

\begin{algorithm}[H]
\caption{$W_1$ empírica en 1D}
\label{alg:ap_w1}
\KwIn{Muestras $\{x_i\}$ de $p$, $\{y_j\}$ de $q$}
\KwOut{$\hat{W}_1(p,q)$}

Ordenar: $x_{(1)}\leq\cdots\leq x_{(N)}$,
$y_{(1)}\leq\cdots\leq y_{(M)}$\;
Interpolar a grilla común $z_1\leq\cdots\leq z_K$ con $K=\max(N,M)$\;
$F_p(z_k) \leftarrow$ CDF empírica de $\{x_i\}$ en $z_k$\;
$F_q(z_k) \leftarrow$ CDF empírica de $\{y_j\}$ en $z_k$\;
$\hat{W}_1 \leftarrow \sum_{k=1}^{K-1}|F_p(z_k)-F_q(z_k)|\cdot(z_{k+1}-z_k)$\;
\Return $\hat{W}_1$\;
\end{algorithm}

\noindent\textbf{Alternativa directa:} en \texttt{scipy},
\texttt{scipy.stats.wasserstein\_distance(x, y)} implementa esto exactamente.

\subsection{KL empírica por estimación con histograma}
\label{subsec:ap_kl_empirica}

Para estimar $D_{\mathrm{KL}}(p\|q)$ a partir de muestras:

\begin{algorithm}[H]
\caption{KL empírica (histograma)}
\label{alg:ap_kl_emp}
\KwIn{Muestras $\{x_i\}$ de $p$, $\{y_j\}$ de $q$; número de bins $B$}
\KwOut{$\hat{D}_{\mathrm{KL}}(p\|q)$}

Calcular rango común $[a,b]$ que contenga todas las muestras\;
Construir histogramas normalizados $\hat{p}_k$, $\hat{q}_k$ sobre $B$ bins en $[a,b]$\;
Sumar $1/N$ (suavizado de Laplace) para evitar bins vacíos: $\hat{p}_k \leftarrow
(\hat{p}_k N + 1)/(N + B)$, ídem para $\hat{q}_k$\;
$\hat{D}_{\mathrm{KL}} \leftarrow \sum_{k=1}^B \hat{p}_k\log(\hat{p}_k/\hat{q}_k)$\;
\Return $\hat{D}_{\mathrm{KL}}$\;
\end{algorithm}

\noindent\textbf{Trampas:}
\begin{itemize}
  \item Sin suavizado, bins con $\hat{q}_k = 0$ producen $\log(0)$ indefinido.
  \item El suavizado de Laplace introduce sesgo que decrece con $N$: para
        $N < 100$ el estimador tiene varianza alta; usar $B \leq \sqrt{N}$ bins.
  \item Para comparaciones entre distribuciones conocidas (gaussiana, OU
        estacionario), calcular la KL analíticamente cuando sea posible en
        lugar de estimarla.
\end{itemize}

\subsection{Estimación de densidad por KDE}
\label{subsec:ap_kde}

Para visualizar la distribución de muestras en los ejercicios de difusión:

\[
  \hat{p}(x) = \frac{1}{Nh}\sum_{i=1}^N K\!\left(\frac{x-x_i}{h}\right),
\]
con kernel gaussiano $K(u) = e^{-u^2/2}/\sqrt{2\pi}$ y ancho de banda óptimo
de Silverman: $h = 1.06\,\hat\sigma\,N^{-1/5}$, donde $\hat\sigma$ es la
desviación estándar de la muestra.

\noindent En \texttt{scipy}: \texttt{scipy.stats.gaussian\_kde(samples)}.
El ancho de banda puede ajustarse con \texttt{bw\_method='silverman'} o
especificando un escalar.

%───────────────────────────────────────────────────────────────────────────────
\section{Denoising score matching: bucle de entrenamiento}
\label{sec:ap_dsm}
%───────────────────────────────────────────────────────────────────────────────

\textit{Usado en: Capítulo~\ref{ch:cap11} (Ejercicio computacional 2),
Capítulo~\ref{ch:cap12} (Ejercicios computacionales 1 y 2).}

\medskip
El bucle de entrenamiento del score network no está escrito de forma compacta
en el cuerpo del texto. Se presenta aquí para los ejercicios que requieren
entrenarlo desde cero.

\begin{algorithm}[H]
\caption{Entrenamiento por denoising score matching (DDPM)}
\label{alg:ap_dsm}
\KwIn{Dataset $\{\bm{x}_0^{(i)}\}$; score network $\bm{\varepsilon}_\theta$;
      schedule $\{\bar\alpha_t\}_{t=1}^T$; optimizador Adam; iteraciones $K$}
\KwOut{Parámetros $\theta$ entrenados}

Inicializar $\theta$, estado de Adam\;
\For{$k = 1, \ldots, K$}{
  $\bm{x}_0 \leftarrow$ muestra aleatoria del dataset\;
  $t \leftarrow \mathcal{U}\{1,\ldots,T\}$\tcp*{nivel de ruido aleatorio}
  $\bm{\varepsilon} \leftarrow \mathcal{N}(\bm{0},\bm{I})$\;
  $\bm{x}_t \leftarrow \sqrt{\bar\alpha_t}\,\bm{x}_0 +
  \sqrt{1-\bar\alpha_t}\,\bm{\varepsilon}$\tcp*{forward de un paso}
  $\hat{\bm{\varepsilon}} \leftarrow \bm{\varepsilon}_\theta(\bm{x}_t, t)$\;
  $\mathcal{L} \leftarrow \|\hat{\bm{\varepsilon}} - \bm{\varepsilon}\|^2$\;
  Paso de Adam sobre $\nabla_\theta\mathcal{L}$\;
}
\end{algorithm}

\noindent\textbf{El tiempo $t$ como entrada a la red.} La red necesita saber
en qué nivel de ruido opera. La forma estándar es inyectar un \textbf{embedding
sinusoidal} del tiempo:
\[
  \bm{e}(t) = \left[\sin\!\left(\frac{t}{10^{4k/d}}\right),
  \cos\!\left(\frac{t}{10^{4k/d}}\right)\right]_{k=0}^{d/2-1}\in\mathbb{R}^d,
\]
que luego se proyecta con una capa densa y se suma a los mapas de características
de la red. Para los ejercicios 1D, concatenar simplemente el escalar $t/T$
normalizado a la entrada es suficiente.

\medskip
\noindent\textbf{Trampas:}
\begin{itemize}
  \item Muestrear $t$ \textbf{uniformemente} sobre $\{1,\ldots,T\}$, no usar
        un valor fijo ni un orden secuencial. El entrenamiento necesita ver todos
        los niveles de ruido por igual.
  \item La pérdida no está ponderada (pesos $\lambda_t = 1$): esto es una
        simplificación válida que empíricamente mejora la calidad perceptual
        pero es una desviación del ELBO exacto del Teorema~\ref{thm:ddpm_elbo}.
  \item Para datos 1D, el score network puede ser tan simple como un MLP de 3
        capas con 128 neuronas. No es necesario un U-Net hasta que la dimensión
        supere $\sim 100$.
\end{itemize}

%───────────────────────────────────────────────────────────────────────────────
\section{Tabla de correspondencia: función $\to$ librería}
\label{sec:ap_librerias}
%───────────────────────────────────────────────────────────────────────────────

\textit{Referencia rápida para los ejercicios computacionales. Los tres
frameworks son equivalentes para los propósitos del libro; la elección es
personal. NumPy/SciPy es el punto de partida natural para físicos; JAX
comparte la API de NumPy y añade diferenciación automática; PyTorch es el
estándar en investigación de ML.}

\medskip
Las listas se organizan por framework. Cada entrada tiene el formato
\emph{operación}: \texttt{llamada}.

\subsubsection*{NumPy / SciPy}
\begin{itemize}[leftmargin=*, nosep]
  \item Pseudoinversa $\bm{A}^\dagger$:
        \texttt{np.linalg.pinv(A)}
  \item SVD $\bm{A}=\bm{U}\bm{\Sigma}\bm{V}^\top$:
        \texttt{np.linalg.svd(A)}
  \item Norma espectral $\|\bm{A}\|_2$:
        \texttt{np.linalg.norm(A, 2)}
  \item Norma de Frobenius $\|\bm{A}\|_F$:
        \texttt{np.linalg.norm(A, 'fro')}
  \item Mínimos cuadrados $\min\|\bm{A}\bm{x}-\bm{b}\|$:
        \texttt{np.linalg.lstsq(A, b)}
  \item Producto acumulado $\bar\alpha_t = \prod_s\alpha_s$:
        \texttt{np.cumprod(alpha)}
  \item Ruido $\bm{\varepsilon}\sim\mathcal{N}(\bm{0},\bm{I})$:
        \texttt{np.random.randn(d)}
  \item Distancia $W_1$ en 1D:
        \texttt{scipy.stats.wasserstein\_distance(x, y)}
  \item KDE gaussiana:
        \texttt{scipy.stats.gaussian\_kde(samples)}
  \item Gradiente por diferencias finitas:
        \texttt{scipy.optimize.approx\_fprime(x, f, eps)}
\end{itemize}

\subsubsection*{PyTorch}
\begin{itemize}[leftmargin=*, nosep]
  \item Pseudoinversa:
        \texttt{torch.linalg.pinv(A)}
  \item SVD:
        \texttt{torch.linalg.svd(A)}
  \item Norma espectral:
        \texttt{torch.linalg.matrix\_norm(A, 2)}
  \item Norma de Frobenius:
        \texttt{torch.linalg.matrix\_norm(A, 'fro')}
  \item Mínimos cuadrados:
        \texttt{torch.linalg.lstsq(A, b)}
  \item Producto acumulado:
        \texttt{torch.cumprod(alpha, dim=0)}
  \item Ruido gaussiano:
        \texttt{torch.randn(d)}
  \item Backward pass (gradiente de la pérdida):
        \texttt{loss.backward()}
  \item Gradiente explícito:
        \texttt{torch.autograd.grad(loss, params)}
  \item Jacobiano:
        \texttt{torch.autograd.functional.jacobian(f, x)}
  \item Hessiano:
        \texttt{torch.autograd.functional.hessian(f, x)}
\end{itemize}

\subsubsection*{JAX}
\begin{itemize}[leftmargin=*, nosep]
  \item Pseudoinversa:
        \texttt{jnp.linalg.pinv(A)}
  \item SVD:
        \texttt{jnp.linalg.svd(A)}
  \item Norma espectral:
        \texttt{jnp.linalg.norm(A, 2)}
  \item Norma de Frobenius:
        \texttt{jnp.linalg.norm(A, 'fro')}
  \item Mínimos cuadrados:
        \texttt{jnp.linalg.lstsq(A, b)}
  \item Producto acumulado:
        \texttt{jnp.cumprod(alpha)}
  \item Ruido gaussiano \emph{(requiere clave explícita)}:
        \texttt{jax.random.normal(key, (d,))}
  \item Gradiente $\nabla_\theta\mathcal{L}$:
        \texttt{jax.grad(loss)(theta)}
  \item Jacobiano:
        \texttt{jax.jacobian(f)(x)}
  \item Hessiano:
        \texttt{jax.hessian(f)(x)}
  \item JVP (modo forward):
        \texttt{jax.jvp(f, (x,), (v,))}
  \item VJP (modo reverse):
        \texttt{jax.vjp(f, x)}
\end{itemize}

\medskip
\noindent\textbf{Nota sobre la aleatoriedad en JAX.} A diferencia de NumPy y
PyTorch, JAX no tiene estado aleatorio global. Cada llamada aleatoria requiere
una clave explícita: \texttt{key = jax.random.PRNGKey(0)}. Para generar
múltiples muestras independientes en un bucle, dividir la clave en cada paso:
\texttt{key, subkey = jax.random.split(key)}.

\medskip
\noindent\textbf{Cuándo usar cada uno.}
Para los ejercicios de SDEs (Caps.~\ref{ch:cap10}--\ref{ch:cap11}),
NumPy es suficiente. Para entrenamiento de redes (Caps.~\ref{ch:cap8},
\ref{ch:cap12}), los tres son equivalentes; PyTorch tiene la ventaja de que
el código de los papers citados (DDPM, score SDE, flow matching) está
publicado en ese framework. JAX es superior para diferenciación de segundo
orden (hessianos, JVPs): su API es más directa y su sistema de
compilación JIT hace los bucles de simulación más rápidos.

%───────────────────────────────────────────────────────────────────────────────
\begin{notasbib}
Los métodos numéricos para SDEs se presentan en detalle en Kloeden y Platen
(1992)~\cite{kloeden1992numerical}: la referencia estándar del campo. El orden
de convergencia fuerte y débil de Euler-Maruyama y Milstein está demostrado en
los Capítulos 9 y 14 de esa referencia. Para una presentación más accesible
orientada a simulaciones estocásticas en física y biología, Higham
(2001)~\cite{higham2001algorithmic} es un tutorial conciso.

El ancho de banda óptimo de Silverman para KDE está demostrado en Silverman
(1986)~\cite{silverman1986density}. La distancia de Wasserstein $W_1$ y su
cálculo por cuantiles en 1D están en Villani (2009)~\cite{villani2009optimal},
Capítulo 2. El estimador de KL por histograma con suavizado de Laplace y sus
propiedades de convergencia se discuten en Poczos y Schneider (2011).

El embedding sinusoidal para el tiempo en los modelos de difusión fue introducido
por Vaswani et al.\ (2017)~\cite{vaswani2017attention} para codificación
posicional y adoptado por Ho et al.\ (2020)~\cite{ho2020denoising} para el
condicionamiento temporal en DDPM.
\end{notasbib}

\printbibliography[heading=subbibliography]